\section{Conclusion and decision use}
Zurich minimizes maximum regret under the retained baseline and equal weights,
while all sixteen cities remain possible champions. The recommendation gives a
family a defensible starting comparison under explicit criteria and source
bounds. Education priority changes that starting choice to Oslo, demonstrating
a practical reason to discuss priorities before using a city order.

Verified supplements reduce average score-interval width by 13.97\% without
replacing the original table's values. The controlled formulation comparison
then holds that evidence fixed: retaining uncertainty changes thirteen decision
ranks and preserves completions in which another city exceeds Zurich. These
are distinct improvements in information and interpretation, rather than a
claim that the model has measured a city's true superiority.

The next evidence priorities follow the unresolved decision. Comparable
high-school travel distributions, matched pupil--teacher populations and
complete city-boundary facility inventories would tighten the main bounds.
Matched student locations, prices, eligibility and opening hours would support
an access analysis. Comparable multi-year observations would allow future
scenarios to be calibrated to data. Until those are available, the city choice,
preference sensitivity and development scenarios retain their stated scopes.
